%=======================================================================

% Template for Thesis at Computer Sciencte at
% Ostfalia University of Applied Science
% Edit: Dirk J. Lehmann, Kai Michael Blum

%=======================================================================
\documentclass[
		11pt,
		a4paper, 
		oneside, % twoside use two sided for book layout
		% openany,
		toc=listofnumbered,
		toc=bibliography, 
		toftofoc,
		headings=small,
		headings=twolineappendix
]{scrbook}
\usepackage[utf8]{inputenc}

%=======================================================================

\newif\ifdraft
\draftfalse % Sagt aus, dass dieses Dokument ein Entwurf ist. Somit wird todonotes aktiviert. Zum deaktivieren diese Zeile auskommentieren oder auf \draftfalse setzen.

%=======================================================================

% In diesem Dokument werden die globalen \usepackage{}-Befehle eingefügt.
% Definition des Seitestils (Ränder, Schriftart, Abstände, usw.)

\usepackage[
		top=2.5cm, 
		bottom=2cm, 
		left=3cm, 
		right=2cm,
		includeheadfoot
]{geometry}

\usepackage[german]{babel}
\selectlanguage{german}
\usepackage[scaled]{Meta/uarial}     % Schriftart Arial
\usepackage[T1]{fontenc}        % bessere und richtige Schriftausgabe

% here for tables
\usepackage{here}

% appendix
\usepackage[toc,page]{appendix}
\renewcommand\appendixpagename{Anhang}
\renewcommand\appendixtocname{Anhang}

% chapter 14pt, section 12pt, subsection 11pt
\setkomafont{chapter}{\normalfont\bfseries\fontsize{14}{16.8}\selectfont}   
\setkomafont{section}{\normalfont\bfseries\fontsize{12}{14.4}\selectfont}   
\setkomafont{subsection}{\normalfont\bfseries\selectfont}     

\usepackage{etoolbox}
\makeatletter
% No breake in chapter
\patchcmd{\@@makechapterhead}{\endgraf\nobreak\vskip.5\baselineskip}{}{}{}


\let\my@chapter\@chapter
\renewcommand*{\@chapter}{%
	\addtocontents{lof}{\protect\addvspace{0pt}}%
	\my@chapter}

\makeatother

%\setuptoc{toc}		% Inhaltsverzeichnis im Inhaltsverzeichnis

\usepackage[titles]{tocloft} % Anhang im Inhaltsverzeichnis

% Abstand Chapter
\renewcommand*{\chapterheadstartvskip}{\vspace*{-\baselineskip}}
\renewcommand*{\chapterheadendvskip}{\vspace*{10pt}}

% durchgehende nummerierung bei Tabellen und Bildern
\usepackage{chngcntr}
\counterwithout{figure}{chapter}
\counterwithout{table}{chapter}

% Prefixe von Abbildungen und Tabellen
\addto\captionsngerman{\renewcommand{\figurename}{Fig.}}
\addto\captionsngerman{\renewcommand{\tablename}{Tab.}}
\renewcommand{\cftfigpresnum}{\figurename\ }
\renewcommand{\cfttabpresnum}{\tablename\ }

% Abstand zwischen Nummer und Name von Abbildungen und Tabellen in Verzeichnissen
\newlength{\mylenf}
\settowidth{\mylenf}{\cftfigpresnum}
\setlength{\cftfignumwidth}{\dimexpr\mylenf+1.5em}
\setlength{\cfttabnumwidth}{\dimexpr\mylenf+1.5em}

% Einrückung in Tabellen- und Abbildungsverzeichnis aufheben
\renewcommand{\cfttabindent}{0em}
\renewcommand{\cftfigindent}{0em}


\usepackage[bottom]{footmisc} % Fußnoten nach unten stellen

\usepackage{amsmath} % amsmath package verwenden
\usepackage{enumitem} % enumaration für listen modifizierbar machen
\usepackage{amsfonts} % amsfonts package verwenden
\usepackage[linesnumbered,ruled]{algorithm2e}
\SetAlgorithmName{Heuristic}{heuristic}{List of Heuristics}

\usepackage[noend]{algpseudocode}
\usepackage{pgfplots}
\pgfplotsset{compat=1.18}
\usepackage{listings}

\setlength{\parindent}{0em} % Absatzeinrückung
\setlength{\parskip}{0.5em}  % Absatzeinrückung

%1,5 facher Zeilenabstand
\usepackage{setspace}
\onehalfspacing			

% Grafiken einbinden
\usepackage{graphicx}

% mainmatter erweitern
% Kopfzeile auf auf seiten mit Kapitelstart
\newcommand{\oldmainmatter}{}
\let\oldmainmatter\mainmatter
\renewcommand{\mainmatter}{\oldmainmatter \renewcommand{\chapterpagestyle}{fancy}}

% Kopf und Fußzeilen
\usepackage{fancyhdr}
\pagestyle{fancy}
\setlength{\headheight}{14.5pt}
\setlength{\headsep}{0.7cm}
\fancyhf{}
\rhead{\nouppercase\rightmark}
\lhead{\nouppercase\leftmark}
\rfoot{Seite \thepage}

% Kapitel

\usepackage{csquotes}
\usepackage[square,numbers]{natbib} % Natbib für DINAT

% PDF Einstellungen
\usepackage[pdftex, pdfborder={0 0 0}]{hyperref}
%\hypersetup{
%	pdfauthor={\documentname\ \documentmanr},
%	pdftitle={\documenttitle}
%	pdfkeywords={Betreuer: \documentbetreuer}
%	}

% Besseres Aussehen in PDF
\usepackage{lmodern}

% Anhang zusätzliche einstellungen:
% Nummerierung ändern:
\usepackage{ifthen}

\newcommand{\oldappendix}{}
\let\oldappendix\appendix
\renewcommand{\appendix}{
	
	\oldappendix
	\newcommand{\oldchapter}{}
	\let\oldchapter\chapter
	\renewcommand{\chapter}[2][]{
		\ifthenelse{\equal{##1}{}}  {
			\oldchapter{##2}
		}{
			\oldchapter[##1]{##2}
		}    		% if title given
		\setcounter{page}{1}
	}    
	    
	\renewcommand{\thepage}{\Alph{chapter} \arabic{page}}
	
	%\cftaddtitleline{toc}{chapter}{\appendixname}{}
	\addtocontents{toc}{\protect\setcounter{tocdepth}{-1}}
	\hypersetup{bookmarksdepth=-1}
	\renewcommand*{\autodot}{:}% Doppelpunkt nur im Anhang
}

\newcommand{\postappendix}{
	\addtocontents{toc}{\protect\setcounter{tocdepth}{3}}
	\hypersetup{bookmarksdepth=3}
}


% Todo package ein- und ausschalten
\ifdraft
	\usepackage{todonotes}
\else
	\usepackage[disable]{todonotes}
\fi

% Abkürzungsverzeichnis
\usepackage[printonlyused]{acronym}

% Ränder testen
% \usepackage{showframe} % Auskommentieren, um die Layoutgrenzen anzuzeigen

% Testtexte
\usepackage{lipsum}


%=======================================================================

\newcommand{\documenttitle}{Infrastruktur-agnostische Entwicklung und Bereitstellung von Webanwendung}
\newcommand{\documentname}{Niklas Röske}
\newcommand{\documentmanr}{70456600}
\newcommand{\documentmodul}{}
\newcommand{\documentstudiengang}{Informatik}
\newcommand{\documentachievement}{zur Erlangung des akademischen Grades:\\Master of Science\\}




\newcommand{\documentbetreuer}{Prof. Dr. Hans Grönniger}
\newcommand{\documentbetreuerzwei}{Prof. Dr. Bernd Müller}
\newcommand{\documentsemester}{2023}

\newcounter{documentfakultaet}
\setcounter{documentfakultaet}{1} % 0 = Bau Wasser Boden, 1 = Informatik

%=======================================================================


\begin{document}	

    % = = = = =	= = = = = = = = = = = = = = = = = = = = = = = =
	
	% Deckblatt

\begin{titlepage}
    \newgeometry{top=2cm, bottom=2.5cm, left=2.5cm, right=2cm}
    
    \ifnum\value{documentfakultaet}=0
    	\definecolor{strich}{cmyk}{0,.97,.92,.11} % Bau-Wasser-Boden
    \fi\ifnum\value{documentfakultaet}=1
    	\definecolor{strich}{rgb}{0,.62,.89} % Informatik
    \fi
    
    \definecolor{text}{cmyk}{1,.55,0,.53} 
    %\hfill
    \begin{flushleft}
    \includegraphics[scale=0.9]{Bilder/ostfalia_logo2.png}\\
    %\vspace{2em}
    \color{strich} 
    \hspace{24mm}\rule{70mm}{0.25mm}\\
    \hspace{24mm}\Large
    \textsf{Fakultät
        \ifnum\value{documentfakultaet}=0
	    	Bau-Wasser-Boden
	    \fi\ifnum\value{documentfakultaet}=1
	    	Informatik
	    \fi
	}
    \end{flushleft}
    
    \begin{flushleft}
         \vspace{3em}
        {\color{text}\fontsize{21}{20}{\textsf{\documenttitle}}}\\
        \vspace{6em}
         {\large\textsf{\textbf{\documentname}}}\\
         {\small\textsf{Matrikel-Nr. \documentmanr}}\\
        \vspace{3em}
        {\color{black}\large\textsf{Masterarbeit im Studiengang \documentstudiengang\\
        \documentachievement
        \ \\
        %im Studiengang \documentstudiengang an der\\
        Ostfalia Hochschule für angewandte Wissenschaften\\
        %– Hochschule Braunschweig/Wolfenbüttel\\
        \color{strich}
        \hrulefill\\
        \ \\\color{black}
        1. Prüfer: \documentbetreuer
        \ \\
        2. Prüfer: \documentbetreuerzwei
        \ \\
        \ \\
        }}
        
    \end{flushleft}
    \begin{flushright}
    %\vspace{14em}
    \vfill
    \color{text}\small\textsf{
		\ifnum\value{documentfakultaet}=0
    		Salzgitter $ \cdot $ \textbf{Suderburg} $ \cdot $ Wolfenbüttel $ \cdot $ Wolfsburg
    	\fi\ifnum\value{documentfakultaet}=1
    		Salzgitter $ \cdot $ Suderburg $ \cdot $ \textbf{Wolfenbüttel} $ \cdot $ Wolfsburg
    	\fi
    } 
    \end{flushright}
    
\end{titlepage}

\restoregeometry

    
    % = = = = =	= = = = = = = = = = = = = = = = = = = = = = = =	
	
	\frontmatter
	\pagenumbering{Roman}
	\ifnum\value{documentfakultaet}=1
		\chapter*{Erklärung}

Hiermit versichere ich, dass ich die vorliegende Arbeit selbständig verfasst und keine anderen als die angegebenen Quellen und Hilfsmittel benutzt habe. Ich versichere, dass ich alle wörtlich oder sinngemäß aus anderen Werken übernommenen Aussagen als solche gekennzeichnet habe, und dass die eingereichte Arbeit weder vollständig noch in wesentlichen Teilen Gegenstand eines anderen Prüfungsverfahrens gewesen ist.

\vspace*{3em}

\newcommand*{\SignatureAndDate}[1]{
	\par\noindent\makebox[52mm]{\hrulefill}     \hfill\makebox[65mm]{\hrulefill}
	\par\noindent\makebox[52mm][l]{Ort, Datum}	\hfill\makebox[62mm][l]{#1}
}

\SignatureAndDate{Unterschrift}
		\markboth{Erklärung}{}
	\fi
    
    % = = = = =	= = = = = = = = = = = = = = = = = = = = = = = =
	
	\chapter*{Abstract}

Your abstract goes here..
    
    % = = = = =	= = = = = = = = = = = = = = = = = = = = = = = =	
	
	% disable header
	\fancyhead[L]{\chaptermark}  % L -> Left part
	\fancyhead[RO]{\chaptermark} % RO -> Right part on Odd pages
	
	\tableofcontents % Inhaltsverzeichnis
	
	% = = = = =	= = = = = = = = = = = = = = = = = = = = = = = =
	
	\chapter{Abkürzungsverzeichnis}
%
\begin{itemize}
    \item
        OSC...Orthographic Star Coordinates
    \item 
        SC...Star Coordinates
    \item 
        CO...Composition Operators
    \item
    	LSS...Least Square Solution
    \item 
        DSC...Distance Consistency
    \item  
        CD...Centroid Density
    \item
        CDC...Centroid Distance Change
    \item
        VML...Visual Machine Learning
\end{itemize}







	
	% = = = = =	= = = = = = = = = = = = = = = = = = = = = = = =
	
	\mainmatter
	
	% reanable header
	\fancyhf{}
	\rhead{\nouppercase\rightmark}
	\lhead{\nouppercase\leftmark}
	
	\fancypagestyle{plain}{}
	%\fancyhead[LE]{\chaptermark} % Chapter in header Left
	%\fancyhead[RE]{\chaptermark} % Page number in header Right
	%\fancyfoot[LE]{\thepage}     % LE -> Left part on Even pages
	\fancyfoot[RO]{\thepage}     % RO -> Right part on Odd pages
	
	% = = = = =	= = = = = = = = = = = = = = = = = = = = = = = =
	
	% reset list of used acronyms
	\acresetall
	
	% = = = = =	= = = = = = = = = = = = = = = = = = = = = = = =
	
	% Kapitel einfügen
	\chapter{Einleitung}

\section{Hintergrund und Motivation}
Selbst in einer schnelllebigen technologischen Landschaft bleibt eine Konstante bestehen: 
die kontinuierliche Evolution von Webanwendungen und ihrer zugrunde liegenden Infrastrukturen. 
Mit der stetigen Weiterentwicklung von Softwareentwicklungsmethoden, Bereitstellungstechnologien 
und Cloud-Plattformen haben sich auch die Ansätze zur Entwicklung und Bereitstellung von Webanwendungen stark verändert. 
In diesem Kontext gewinnt das Konzept der infrastruktur-agnostischen Entwicklung und Bereitstellung zunehmend an Bedeutung. 
Infrastruktur-agnostische Ansätze ermöglichen es, Webanwendungen unabhängig von den spezifischen technologischen Details der zugrunde liegenden Infrastruktur zu gestalten und bereitzustellen. 
Dies eröffnet Chancen für eine erhöhte Flexibilität, Skalierbarkeit und Effizienz in der Softwareentwicklung und -bereitstellung.


\section{Zielsetzung der Arbeit}
In dieser Masterarbeit wird die Thematik der infrastruktur-agnostischen Entwicklung und Bereitstellung von Webanwendungen umfassend untersucht. 
Dabei werden die zugrunde liegenden Konzepte, Methoden und Technologien analysiert, die es ermöglichen, Webanwendungen von den Details der Infrastruktur zu abstrahieren. 
Ein besonderer Fokus liegt dabei auf der Identifizierung von Best Practices, Herausforderungen und potenziellen Lösungsansätzen im Zusammenhang mit der Umsetzung infrastruktur-agnostischer Ansätze. 
Des Weiteren werden Fallstudien und Praxisbeispiele aus verschiedenen Industriezweigen betrachtet, um ein umfassendes Verständnis für die Anwendung und Auswirkungen dieser Ansätze zu entwickeln.

Indem diese Arbeit die Vor- und Nachteile, die möglichen Auswirkungen auf die Softwareentwicklung und die langfristigen Perspektiven der infrastruktur-agnostischen Entwicklung und Bereitstellung von Webanwendungen beleuchtet, trägt sie zur Erweiterung des Wissens in diesem sich rasch entwickelnden Bereich bei. 
Letztendlich wird diese Arbeit dazu beitragen, Entwicklern, Unternehmen und Entscheidungsträgern dabei zu helfen, fundierte Entscheidungen im Hinblick auf die Wahl geeigneter Ansätze für die Entwicklung und Bereitstellung ihrer Webanwendungen zu treffen, während sie von den Vorteilen einer infrastruktur-agnostischen Herangehensweise profitieren.

	% ----------------------------------------------------------
% ----------------------------------------------------------
\chapter{Theoretischer Hintergrund}

\section{Web-Anwendungen}
\label{sec:visual_analytics}

\section{Infrastruktur unabhängige Entwicklung}

\section{Herausforderungen bei der Entwicklung}

\cite{Sammon.1969}\\
\cite{Demartines.1997}\\
	\chapter{Methodik}

\section{Auswahl Entwicklungstechnologien}

\section{Architektur und Design der Web-Anwendung}

\section{Implementierung un Testing?}
	\chapter{Empfehlungen und bewährte Praktiken}
\label{sec:approach}

\section{Best Practices}

\section{Technische Aspekte}
	\chapter{Fallstudie}


\section{Beschreibung der Anwendung}

\section{Bewertung der Ergebnisse}

	\chapter{Diskussion}

\section{Wichtigste Erkenntnisse}

\section{Beantwortung der Forschungsfrage}

\section{Kritische Bewertung der EMpfehlung}
	\chapter{Ausblick}

Your discussion goes here ...

    % = = = = =	= = = = = = = = = = = = = = = = = = = = = = = =

    % Bilbliography
	\bibliographystyle{dinat} 
	\bibliography{literatur}
	
	% = = = = =	= = = = = = = = = = = = = = = = = = = = = = = =
	
	\ifnum\value{documentfakultaet}=0
		\input{erklaerung}
		\markboth{Erklärung}{}
	\fi
	
	% = = = = =	= = = = = = = = = = = = = = = = = = = = = = = =
	
    %\addtocontents{toc}{\cftpagenumbersoff{chapter}} % Keine Seitenzahl im Inhaltsverzeichnis
	\begin{appendices}              % Anhang Titel-Seite
	\end{appendices}
	
	\fancyfoot[R] {\thepage}        % Seitenzahl Right
	
	\appendix
	\chapter{Pseudocode for Heuristics}
\label{appendix:pseudocode_heuristics}

\begin{figure}[H]
\begin{algorithm}[H]
\DontPrintSemicolon
\SetKwProg{Fn}{function}{:}{}
\Fn{heuristic\_random\_selection\_shift($df\_p, iterator, num\_classes$)}{
\KwData{$df\_p$ is the dataframe in projection space, $iterator$ is the current iteration step, $num\_classes$ is the amount of different classes in the dataset}
\KwResult{$dp$ shift vector, $selected\_class$ to be shifted, calculated $DSC\_value$, calculated $CD\_value$, calculated $total\_dist$}
\tcc{extract star coords and class of data to a new dataframe}
df\_star = df\_p[['X','Y','class']]\;
\tcc{calculate class centroids and save them in a new dataframe}
df\_centroids = df\_star.groupby('class', sort=True).mean().reset\_index()\;
\tcc{calculate coordinates of the central centroid}
central\_centroid = df\_centroids[['X', 'Y']].mean()\;
\tcc{call a function to calculate all distances between centroids}
df\_distances = calc\_centroid\_distances(df\_centroids)\;
\tcc{calculate the distance from each point to its associated centroid}
df\_centroid\_distances = calc\_dist\_p\_to\_assoc\_centroid(df\_centroids, df\_star)\;
\tcc{calculate CD, DSC and total\_dist for result}
CD\_value = df\_centroid\_distances['distance'].sum()\;
DSC\_value = calc\_dsc(df\_centroids, df\_star)\;
total\_dist = df\_distances['distance'].sum()\;
\color{green}
\tcc{randomly choose a class to be shifted}
selected\_class = np.random.randint(num\_classes)\;
\color{black}
\tcc{select all distances for selected class}
selected\_class\_distances = df\_distances.(df\_distances[selected\_class])\;
\tcc{other\_class is the nearest class to selected\_class}
min\_selected\_class\_distance = selected\_class\_distances['distance'].idxmin()\;
\tcc{calculate the new shifting vector dp}
dp = calc\_dp(df\_centroids, selected\_class, other\_class, central\_centroid, num\_iter)\;
}
\caption{Random Selection Shift RSS}
\label{func:rss}
\end{algorithm}
\caption{Pseudocode for RSS}
\end{figure}

\begin{figure}[H]
\begin{algorithm}[H]
\DontPrintSemicolon
\SetKwProg{Fn}{function}{:}{}
\Fn{heuristic\_order\_selection\_shift($df\_p, iterator, num\_classes$)}{
\KwData{$df\_p$ is the dataframe in projection space, $iterator$ is the current iteration step, $num\_classes$ is the amount of different classes in the dataset}
\KwResult{$dp$ shift vector, $selected\_class$ to be shifted, calculated $DSC\_value$, calculated $CD\_value$, calculated $total\_dist$}
\tcc{extract star coords and class of data to a new dataframe}
df\_star = df\_p[['X','Y','class']]\;
\tcc{calculate class centroids and save them in a new dataframe}
df\_centroids = df\_star.groupby('class', sort=True).mean().reset\_index()\;
\tcc{calculate coordinates of the central centroid}
central\_centroid = df\_centroids[['X', 'Y']].mean()\;
\tcc{call a function to calculate all distances between centroids}
df\_distances = calc\_centroid\_distances(df\_centroids)\;
\tcc{calculate the distance from each point to its associated centroid}
df\_centroid\_distances = calc\_dist\_p\_to\_assoc\_centroid(df\_centroids, df\_star)\;
\tcc{calculate CD, DSC and total\_dist for result}
CD\_value = df\_centroid\_distances['distance'].sum()\;
DSC\_value = calc\_dsc(df\_centroids, df\_star)\;
total\_dist = df\_distances['distance'].sum()\;
\color{green}
\tcc{choose a class to be shifted by order}
selected\_class = num\_iter \% num\_classes)\;
\color{black}
\tcc{select all distances for selected class}
selected\_class\_distances = df\_distances.(df\_distances[selected\_class])\;
\tcc{other\_class is the nearest class to selected\_class}
min\_selected\_class\_distance = selected\_class\_distances['distance'].idxmin()\;
\tcc{calculate the new shifting vector dp}
dp = calc\_dp(df\_centroids, selected\_class, other\_class, central\_centroid, num\_iter)\;
}
\caption{Order Selection Shift - OSS}
\label{func:oss}
\end{algorithm}
\caption{Pseudocode for OSS}
\end{figure}

\begin{figure}[H]
\begin{algorithm}[H]
\DontPrintSemicolon
\SetKwProg{Fn}{function}{:}{}
\Fn{heuristic\_order\_selection\_shift($df\_p, iterator, num\_classes$)}{
\KwData{$df\_p$ is the dataframe in projection space, $iterator$ is the current iteration step, $num\_classes$ is the amount of different classes in the dataset}
\KwResult{$dp$ shift vector, $selected\_class$ to be shifted, calculated $DSC\_value$, calculated $CD\_value$, calculated $total\_dist$}
\tcc{extract star coords and class of data to a new dataframe}
df\_star = df\_p[['X','Y','class']]\;
\tcc{calculate class centroids and save them in a new dataframe}
df\_centroids = df\_star.groupby('class', sort=True).mean().reset\_index()\;
\tcc{calculate coordinates of the central centroid}
central\_centroid = df\_centroids[['X', 'Y']].mean()\;
\tcc{call a function to calculate all distances between centroids}
df\_distances = calc\_centroid\_distances(df\_centroids)\;
\tcc{calculate the distance from each point to its associated centroid}
df\_centroid\_distances = calc\_dist\_p\_to\_assoc\_centroid(df\_centroids, df\_star)\;
\tcc{calculate CD, DSC and total\_dist for result}
CD\_value = df\_centroid\_distances['distance'].sum()\;
DSC\_value = calc\_dsc(df\_centroids, df\_star)\;
total\_dist = df\_distances['distance'].sum()\;
\color{green}
\tcc{find the maxmimum distance}
max\_dist\_idx = df\_centroid\_distances['distance'].idxmax()\;
\tcc{select the point that is to be shifted}
centroid\_id = df\_centroid\_distances.loc[max\_dist\_idx, 'class']\;
centroid\_coords = [df\_centroids.loc[centroid\_id, 'X'], df\_centroids.loc[centroid\_id, 'Y']]\;
point\_coord = [df\_centroid\_distances.loc[max\_dist\_idx, 'X'], df\_centroid\_distances.loc[max\_dist\_idx, 'Y']]\;
\tcc{create noise}
noise = np.random.normal(0, 1, 1) * 100 / (1000 + num\_iter)\;
\tcc{calculate the new shifting vector dp}
dp = [centroid\_coords[0] - point\_coord[0] + noise, centroid\_coords[1] - point\_coord[1] + noise]\;
}
\caption{Point Selection Shift - PSS}
\label{func:pss}
\end{algorithm}
\caption{Pseudocode for PSS}
\end{figure}

\chapter{Key values for all heuristics for each dataset}
\label{appendix:key_values_heuristics}

\begin{table}[H]
\centering
    \fbox{%
    \begin{tabular}{ c || c | c | c | c | c | c }
    dataset & $DSC_{start}$ & $DSC_{end}$ & $CD_{start}$ & $CD_{end}$ & $d_{c,total,start}$ & $d_{c,total,end}$ \\
    \hline 
    \hline
    ecoli & 63.88\% & 65.07\% & 11636 & 11651 & 3073 & 3086\\
    \hline
    iris & 89.93\% & 89.93\% & 3585 & 3498 & 340 & 358\\
    \hline
    statlog & 21.06\% & 28.76\% & 65743 & 65963 & 112 & 222\\
    \hline
    wdbc & 86.27\% & 94.72\% & 22820 & 16437 & 79 & 66\\
    \hline
    wine & 72.32\% & 92.66\% & 7703 & 7539 & 201 & 353\\
    \hline
    yeast & 27.44\% & 27.38\% & 44547 & 44911 & 2613 & 2635\\
    \end{tabular}
    }%
    \caption{Key values for RSS for each dataset}
    \label{tab:standard_comparison_var1}
\end{table}

\begin{table}[H]
\centering
    \fbox{%
    \begin{tabular}{ c || c | c | c | c | c | c }
    dataset & $DSC_{start}$ & $DSC_{end}$ & $CD_{start}$ & $CD_{end}$ & $d_{c,total,start}$ & $d_{c,total,end}$ \\
    \hline 
    \hline
    ecoli & 63.88\% & 63.88\% & 11636 & 11652 & 3073 & 3078\\
    \hline
    iris & 89.93\% & 90.6\% & 3585 & 3483 & 340 & 364\\
    \hline
    statlog & 21.06\% & 22.76\% & 65743 & 65784 & 112 & 132\\
    \hline
    wdbc & 86.27\% & 94.72\% & 22820 & 16399 & 79 & 65\\
    \hline
    wine & 72.32\% & 93.22\% & 7703 & 7549 & 201 & 357\\
    \hline
    yeast & 27.44\% & 27.44\% & 44547 & 44713 & 2613 & 2616\\
    \end{tabular}
    }%
    \caption{Key values for OSS for each dataset}
    \label{tab:standard_comparison_var2}
\end{table}

\begin{table}[H]
\centering
    \fbox{%
    \begin{tabular}{ c || c | c | c | c | c | c }
    dataset & $DSC_{start}$ & $DSC_{end}$ & $CD_{start}$ & $CD_{end}$ & $d_{c,total,start}$ & $d_{c,total,end}$ \\
    \hline 
    \hline
    ecoli & 63.88\% & 67.76\% & 11636 & 11545 & 3073 & 3136\\
    \hline
    iris & 89.93\% & 93.29\% & 3585 & 3111 & 340 & 444\\
    \hline
    statlog & 21.06\% & 41.27\% & 65743 & 62631 & 112 & 415\\
    \hline
    wdbc & 86.27\% & 96.3\% & 22820 & 18776 & 79 & 98\\
    \hline
    wine & 72.32\% & 93.22\% & 7703 & 7389 & 201 & 384\\
    \hline
    yeast & 27.44\% & 28.32\% & 44547 & 46156 & 2613 & 2765\\
    \end{tabular}
    }%
    \caption{Key values for MSS for each dataset}
    \label{tab:standard_comparison_var3}
\end{table}

\begin{table}[H]
\centering
    \fbox{%
    \begin{tabular}{ c || c | c | c | c }
    dataset & $DSC_{start}$ & $DSC_{end}$ & $CD_{start}$ & $CD_{end}$ \\
    \hline 
    \hline
    ecoli & 63.88\% & 63.39\% & 11636 & 11729\\
    \hline
    iris & 89.93\% & 89.93\% & 3585 & 3585\\
    \hline
    statlog & 21.06\% & 22.51\% & 65743 & 57698\\
    \hline
    wdbc & 86.27\% & 91.37\% & 22820 & 17566\\
    \hline
    wine & 72.32\% & 100\% & 7703 & 5641\\
    \hline
    yeast & 27.44\% & 24.81\% & 44547 & 43926\\
    \end{tabular}
    }%
    \caption{Key values for PSS for each dataset}
    \label{tab:standard_comparison_var4}
\end{table}
	%\postappendix
	
	% = = = = =	= = = = = = = = = = = = = = = = = = = = = = = =
	
\end{document}